\documentclass[conference]{IEEEtran}

\usepackage{cite}
\usepackage{amsmath,amssymb}
\usepackage{graphicx}
\usepackage{booktabs}
\usepackage[colorlinks=true,allcolors=blue]{hyperref}

\title{Robust Prompt Routing for Multi-Agent Research Assistants}

\author{\IEEEauthorblockN{Anonymous Research Team}}

\begin{document}

\maketitle

\begin{abstract}
This IEEE-style sample illustrates how the \texttt{codex-latex} skill produces structured technical drafts.
The workflow prioritizes deterministic reference handling, concise method descriptions, and reproducible section ordering.
We show a routing baseline for multi-agent research assistants and summarize expected reporting style.
\end{abstract}

\section{Introduction}
Multi-agent systems require clear routing rules to map user intent to specialized workers.
In manuscript preparation, unstable routing often causes inconsistent section depth, missing references, and higher editorial cost.
This sample demonstrates a robust writing-oriented routing setup.

\section{Related Work}
Classic LaTeX guidance emphasizes logical structure and consistent typesetting conventions \cite{lamport1994}.
Recent agentic workflows extend this with deterministic validation passes before publication.

\section{Method}
We define three routing stages: intent classification, resource selection, and verification.
Given a task $x$, a controller selects an action path $a$ that minimizes expected correction cost:
\begin{equation}
\mathcal{L}(a \mid x) = \alpha E_{\text{format}} + \beta E_{\text{reference}} + \gamma E_{\text{compile}}.
\end{equation}

\section{Experiments}
Table~\ref{tab:ieee-results} reports an illustrative comparison.

\begin{table}[t]
\caption{Illustrative evaluation for routing policies.}
\label{tab:ieee-results}
\centering
\begin{tabular}{lcc}
\toprule
Policy & Compile Success & Rework Time \\
\midrule
Heuristic only & 82\% & 37 min \\
Structured routing & 97\% & 14 min \\
\bottomrule
\end{tabular}
\end{table}

\section{Conclusion}
The sample highlights IEEE-compatible structure and reporting style.
It can be reused as a starter document when sharing the skill with teams that submit to IEEE venues.

\bibliographystyle{IEEEtran}
\bibliography{references}

\end{document}
